\subsubsection{Why Simon's Algorithm Works}\label{sec:why-simons-algorithm-works}

In the previous section, we established the circuit and arrived at:
\begin{equation*}
    \ket{\psi_{3}} = \frac{1}{2^{n}} \sum_{x \in \{0,1\}^{n}} \sum_{y \in \{0,1\}^{n}} (-1)^{x \cdot y} \ket{y} \ket{f(x)}
\end{equation*}
We also said that the purpose of those Hadamards is to create \textbf{interference} between inputs related by the hidden period. Now we can prove exactly how that happens. Simon turns the general phenomenon of interference into a very precise filter.

\highspace
\begin{flushleft}
    \textcolor{Green3}{\faIcon{tools} \textbf{Reorganizing the Sum Using Simon's Promise}}
\end{flushleft}
After the final Hadamards, the state is:
\begin{equation*}
    \ket{\psi_3}
    =
    \frac{1}{2^n}
    \sum_{x\in\{0,1\}^n}
    \sum_{y\in\{0,1\}^n}
    (-1)^{x \cdot y}
    \ket{y} \ket{f(x)}
\end{equation*}
At first sight, there is no explicit term involving $x\oplus a$. However, the sum over $x$ iterates over \textbf{every possible $n$-bit string}. Since $a \in \left\{0, 1\right\}^{n}$, for every $x$ in the sum:
\begin{equation*}
    x \oplus a \in \left\{0, 1\right\}^{n}
\end{equation*}
Therefore, if the sum contains the contribution indexed by $x$:
\begin{equation*}
    (-1)^{x \cdot y} \ket{y} \ket{f(x)}
\end{equation*}
Then somewhere else in the same sum it necessarily also contains the contribution indexed by $x\oplus a$:
\begin{equation*}
    (-1)^{(x\oplus a) \cdot y} \ket{y} \ket{f(x\oplus a)}
\end{equation*}
The transformation $x \mapsto x \oplus a$ does not introduce new elements: it simply \textbf{permutes the elements of} $\{0,1\}^{n}$. Thus, we can reorganize the original sum into pairs $x$ and $x \oplus a$. Now Simon's promise becomes useful. It guarantees that:
\begin{equation*}
    f(x) = f(x \oplus a)
\end{equation*}
Hence the two terms of each pair contributes to the \textbf{same final basis state} $\ket{y} \ket{f(x)}$. We can therefore combine their amplitudes:
\begin{gather*}
    (-1)^{x \cdot y} \ket{y} \ket{f(x)} + (-1)^{(x\oplus a) \cdot y} \ket{y} \ket{f(x\oplus a)} \\
    = \left[(-1)^{x \cdot y} + (-1)^{(x\oplus a) \cdot y}\right] \ket{y} \ket{f(x)}
\end{gather*}
Thus, the combined amplitude associated with a Simon pair is proportional to:
\begin{equation*}
    (-1)^{x \cdot y} + (-1)^{(x\oplus a) \cdot y}
\end{equation*}
Now the question is: \emph{\textbf{do these two amplitudes reinforce each other, or cancel each other?}} That depends entirely on $a \cdot y$.

\newpage

\begin{flushleft}
    \textcolor{Green3}{\faIcon{check-circle} \textbf{Simplify $\left(x \oplus a\right) \cdot y$}}
\end{flushleft}
We want to simplify the expression $(x \oplus a) \cdot y$ because our actual \hl{goal} is \textbf{not} to understand $x$ or $y$, but rather to \hl{extract information about the unknown secret $a$.} And simplifying $(x \oplus a) \cdot y$ will allow us to see how $a$ affects the relative sign of the two amplitudes.

\highspace
At this point, we have obtained the combined amplitude:
\begin{equation*}
    \left(-1\right)^{x \cdot y} + \left(-1\right)^{(x\oplus a) \cdot y}
\end{equation*}
And we want to answer the question \emph{for a given $y$, do these two amplitudes add or cancel?} Right now that's difficult to see because the second exponent mixes everything together: $(x \oplus a) \cdot y$. We want to expose explicitly \textbf{how $a$ affects the relative sign} of the two terms.

\highspace
So we rewrite the second exponent $(-1)^{(x \oplus a) \cdot y}$ using arithmetic over $\mathbb{Z}_{2}$ (i.e., modulo 2):
\begin{align*}
    (x \oplus a) \cdot y
    &= \sum_{i=1}^{n} (x_i \oplus a_i)y_i \pmod 2 \\[.3em]
    &= \sum_{i=1}^{n} (x_i + a_i)y_i \pmod 2 \\[.3em]
    &= x \cdot y + a \cdot y \pmod 2
\end{align*}

\begin{remarkbox}[: $\oplus$ vs $+$]
    The equality:
    \begin{equation*}
        x_i \oplus a_i \equiv x_i + a_i \pmod 2
    \end{equation*}
    Holds in \textbf{arithmetic over $\mathbb{Z}_2$}. Indeed, let's check the truth table:
    \begin{center}
        \begin{tabular}{@{} c|c|c|c @{}}
            \toprule
            $x_i$ & $a_i$ & $x_i \oplus a_i$ & $x_i + a_i \pmod 2$ \\ \midrule
            0 & 0 & 0 & 0 \\
            0 & 1 & 1 & 1 \\
            1 & 0 & 1 & 1 \\
            1 & 1 & 0 & 0 \\
            \bottomrule
        \end{tabular}
    \end{center}
    Thus, $1\oplus1=0$, while ordinary integer addition gives $1+1=2$. Reducing modulo $2$ makes them equal.
\end{remarkbox}

\noindent
Therefore, we can rewrite the combined amplitude as:
\begin{align*}
    (-1)^{x \cdot y} + (-1)^{(x\oplus a) \cdot y}
    &= (-1)^{x \cdot y} + (-1)^{x \cdot y + a \cdot y} \\[.3em]
    &\downarrow \text{since } (-1)^{p+q} = (-1)^p(-1)^q \text{ for any integers } p,q \\[.3em]
    &= (-1)^{x \cdot y} + (-1)^{x \cdot y}(-1)^{a \cdot y} \\[.3em]
    &= (-1)^{x \cdot y} \left[1 + (-1)^{a \cdot y}\right]
\end{align*}
This is the expression we wanted! Because now something very interesting happens. Since we're working modulo 2, $a \cdot y$ can only be $0$ or $1$. Therefore, there are only \textbf{two cases}:
\begin{itemize}
    \item $a \cdot y = 0$, then:
        \begin{equation*}
            (-1)^{a \cdot y} = (-1)^0 = 1 \implies 1 + (-1)^{a \cdot y} = 1 + 1 = 2
        \end{equation*}
        Our combined amplitude becomes $2(-1)^{x \cdot y}$, which is \textbf{non-zero}:
        \begin{equation*}
            (-1)^{x \cdot y} \left[1 + (-1)^{a \cdot y}\right] = (-1)^{x \cdot y} \cdot 2 = 2(-1)^{x \cdot y}
        \end{equation*}
        Therefore, the two contributions reinforce each other (\textbf{constructive interference}), and the corresponding $y$ is potentially \textbf{measurable}.

    \item $a \cdot y = 1$, then:
        \begin{equation*}
            (-1)^{a \cdot y} = (-1)^1 = -1 \implies 1 + (-1)^{a \cdot y} = 1 - 1 = 0
        \end{equation*}
        The entire combined amplitude becomes $0$:
        \begin{equation*}
            (-1)^{x \cdot y} \left[1 + (-1)^{a \cdot y}\right] = (-1)^{x \cdot y} \cdot 0 = 0
        \end{equation*}
        The two contributions cancel each other (\textbf{destructive interference}), and the corresponding $y$ is \textbf{not measurable}.
\end{itemize}
We have therefore shown that:
\begin{itemize}
    \item $a \cdot y = 0 \implies$ constructive interference.
    \item $a \cdot y = 1 \implies$ destructive interference.
\end{itemize}
And therefore:
\begin{equation*}
    P(y) > 0 \iff a \cdot y = 0 \pmod 2
\end{equation*}
Or in words, \hl{after the final Hadamards, Simon's interference removes every $y$ that is not orthogonal to the hidden string $a$ modulo $2$.}

\highspace
\textcolor{Green3}{\faIcon{question-circle} \textbf{Why are all valid $y$'s equally likely?}} This is one another interesting point. The \hl{interference condition doesn't prefer one valid $y$ over another!} For every $y$ satisfying constructive interference:
\begin{equation*}
    a \cdot y = 0 \pmod 2
\end{equation*}
The factor:
\begin{equation*}
    1 + (-1)^{a \cdot y}
\end{equation*}
Has the same magnitude, $2$. Therefore \hl{all surviving $y$'s have equal probability.} For a non-zero $a$, exactly half of all $2^n$ bitstrings satisfy $a \cdot y = 0$, so there are $2^{n-1}$ valid $y$'s, each with probability:
\begin{equation*}
    P(y) = \frac{1}{2^{n-1}}
\end{equation*}

\begin{examplebox}[: Constructive and Destructive Interference]
    We have $a=01$. For $n=2$, there are four possible measurement outcomes:
    \begin{equation*}
        y \in \{00, 01, 10, 11\}
    \end{equation*}
    Let's calculate $a \cdot y$ for each of them:
    \begin{itemize}
        \item $y=00$: $a \cdot y = 0\cdot 0 + 1 \cdot 0 = 0$ (constructive interference)
        \item $y=01$: $a \cdot y = 0\cdot 0 + 1 \cdot 1 = 1$ (destructive interference)
        \item $y=10$: $a \cdot y = 0\cdot 1 + 1 \cdot 0 = 0$ (constructive interference)
        \item $y=11$: $a \cdot y = 0\cdot 1 + 1 \cdot 1 = 1$ (destructive interference)
    \end{itemize}
    Therefore, the only valid measurement outcomes are $y=00$ and $y=10$, each with probability:
    \begin{equation*}
        P(y=00) = P(y=10) = \frac{1}{2^{2-1}} = \frac{1}{2}
    \end{equation*}
\end{examplebox}

\highspace
\begin{flushleft}
    \textcolor{Green3}{\faIcon{question-circle} \textbf{What did the quantum circuit actually accomplish?}}
\end{flushleft}
Before the final Hadamards, the secret was encoded indirectly as:
\begin{equation*}
    f(x) = f(x \oplus a)
\end{equation*}
That isn't something we can conveniently read from a single measurement. The final Hadamards cause the paired inputs to interfere:
\begin{equation*}
    (-1)^{x \cdot y} + (-1)^{(x\oplus a) \cdot y}
    =
    (-1)^{x \cdot y} \left[1 + (-1)^{a \cdot y}\right]
\end{equation*}
The interference converts that hidden pairing into a completely different kind of information:
\begin{equation*}
    f(x) = f(x \oplus a) \xrightarrow{H^{\otimes n} + \text{ inference}} a \cdot y = 0 \pmod 2
\end{equation*}
So one execution of the quantum circuit effectively transforms \textbf{hidden periodic structure} into \textbf{one linear constraint on the hidden string}. However, \hl{we still haven't found $a$.} We've only measured one $y$ satisfying $a \cdot y = 0$. The natural next question is: \emph{how many such $y$'s do we need, and how do we solve the measuring equations to reconstruct $a$?} We'll answer that in the next section.